\section{Sprint 6 -- External Team Testing}
%=============================================================================

\begin{priorities}
\begin{enumerate}[leftmargin=*, itemsep=1pt, topsep=0pt]
    \item Hand off a clean, well-labeled prototype with a complete test plan.
    \item Observe the ``Silent Observer'' rule: do not coach the testing team.
    \item When \textit{you} are the testing team, be thorough and professional.
\end{enumerate}
\end{priorities}

\begin{faqbox}[Q: What happens during external team testing?]
Another team will receive your prototype and your test plan, and they will execute your tests and document the results. You will \textit{not} be present to explain anything. This is the ultimate test of your documentation quality. Every ambiguity, missing step, or unclear pass/fail criterion in your test plan will surface here.
\end{faqbox}

\begin{tipbox}
\textbf{Prepare a handoff package, not just a prototype.} Before handing your system to the external team, prepare: (1)~a one-page quick-start guide with numbered steps and photos showing how to power on, configure, and operate the system, (2)~all required test equipment, cables, and fixtures laid out and labeled, (3)~clear photographs of the correct system configuration for each test, (4)~any safety precautions highlighted prominently, and (5)~your test plan reviewed one final time as if you have never seen the project. External testing typically reveals three categories of problems: procedure ambiguities (``connect the sensor'' without specifying which port), missing preconditions (the system needs 5 minutes of warm-up but the procedure says nothing about it), and actual system defects that your team never encountered because they operate the system differently.
\end{tipbox}

\begin{checklistbox}[External Testing Handoff Package Checklist]
\begin{enumerate}[leftmargin=*, itemsep=1pt, topsep=0pt]
    \item One-page quick-start guide with numbered steps and annotated photos
    \item Complete test plan with pass/fail criteria for each test (no ambiguity)
    \item All required test equipment, cables, adapters, and fixtures, labeled
    \item Safety precautions sheet (highlighted, not buried in the manual)
    \item Photographs of correct system configuration for each test
    \item Power supply/charger (fully charged batteries if applicable)
    \item Spare consumables (e.g., extra sensor probes, calibration samples)
    \item Data recording forms or templates for the testing team to fill in
    \item Contact information for emergencies only (not for coaching)
\end{enumerate}
\end{checklistbox}

\begin{warnbox}
\textbf{The ``Silent Observer'' Rule:} During external testing, the designing team must remain \textbf{completely silent}. Do not coach, correct, hint, or demonstrate. If the testing team is struggling with a step, that is data; it means your documentation needs improvement. The moment someone says ``Oh, you just need to press it like this,'' you have invalidated the test of your documentation. Resist the urge. Let the test plan speak for itself. This is one of the most valuable learning experiences in the entire course.
\end{warnbox}

\begin{faqbox}[Q: What if the external team breaks something or does a test wrong?]
This is a real-world scenario: products must survive being used by people who did not design them. If a test is performed incorrectly, it likely means your test procedure was unclear. Use this feedback constructively. If something physically breaks, document it and address it in Sprint~7.
\end{faqbox}

\begin{protip}
When you are the \textit{testing} team for another group, be thorough and fair. Your primary goal is not to make the other team look bad; it is to provide honest, constructive feedback that helps them improve. Document everything that is ambiguous or unclear in the test plan; that feedback is a gift to the design team. The quality of your external testing report reflects your own engineering professionalism, and your teammates will return the favor. This is also an opportunity to see how another team approached a different design problem; you will learn things that make you a better engineer.
\end{protip}

%=============================================================================
